La mise en place de ces quatre méthodes focométriques
a permis d'accomplir les \textbf{deux objectifs majeurs} fixés pour ce TP,
à savoir : 

\begin{itemize}
    \item Retrouver la \textbf{distance focale} d'une lentille convergente et d'une lentille divergente,
    respectivement de $+10$ cm et $-10$ cm ;
    \item \textbf{Confronter différentes méthodes} :
    celles se reposant sur les graduations du banc optique (i.e. la \textbf{méthode de Badal})
    se révèlent \textbf{plus rigoureuses} que celles reposant sur la mesure d'images formées sur un écran.
\end{itemize}

En conclusion, la compréhension des théories focométriques et leur mise en pratique
permettent d'obtenir des \textbf{mesures expérimentales précises} en accord avec la réalité.
Néanmoins, la \textbf{confrontation des incertitudes} laisse entendre que ces quatre expériences ne se valent pas :
les méthodes de \textbf{Silbermann} et des \textbf{points conjugués} présentent des incertitudes
nettement \textbf{plus importantes} que les méthodes de \textbf{Badal} et de \textbf{Bessel}.
Toutefois, l'ensemble des \textbf{modèles théoriques} a pu être \textbf{validé} lors de cette séance de TP.